\begin{theorem}[Large-order good-hierarchy elementarity]
Fix a finite loopless multigraph \(G\) with its individual edge identities,
a palette of \(k\) colors, and a distinguished edge \(e_0\).  Let
\(\mathcal E\) be a positive-rung
\noderef{ExtendedTashkinovTree} on an exact terminal tree-sequence \(T\),
constructed from an arbitrary certified input coloring and having generating
coloring \(\varphi_n\).  Suppose that \(\mathcal E\) has
\noderef{MaximumProperty}.  Fix a \noderef{HierarchyBase} \(\mathcal B\),
an actual \noderef{TashkinovHierarchy} \(H\) over \(\mathcal B\), and a
\noderef{GoodHierarchy} certificate for \(H\).

Assume the following full form of (5.1).  For every positive-rung maximum
ETT \(\mathcal L\) on the same \(G,k,e_0\), with arbitrary certified input
and generating colorings and with rung number \(m<n\), all six assertions
below hold.
\begin{enumerate}
\item Its exact terminal vertex set is elementary under its generating
coloring.
\item If \(L_{m+1}\) is its final ladder output, then under every proper
coloring \(\sigma\) leaving exactly \(e_0\) uncolored and stable on
\(L_{m+1}\) relative to the raw stored colors at slot \(m\) and the slot
\(m\) coloring, every pair \(\alpha,\beta\) for which at least one color is
missing on \(L_{m+1}\) is interchangeable on \(L_{m+1}\).
\item At every index \(a<m\) for which slot \(a+1\) is PE with support
\(v\), whenever \(v\) belongs to the exact terminal tree and has
first-appearance index \(j+1\), every proper exactly-one-uncolored coloring
stable on the ladder tree at index \(a+1\), relative to the stored set and
coloring at slot \(a+1\), is stable before \(v\) in the terminal tree
relative to the stored set and coloring at slot \(j\); if \(j>0\), it is
additionally stable on the predecessor ladder tree at slot \(j\).
Moreover, for distinct indices \(a,b<m\) with distinct PE supports
\(v_a,v_b\) at slots \(a+1,b+1\), if both supports belong to the terminal
tree and have the same first-appearance index, then the two raw stored sets
at \(a+1\) and \(b+1\) are disjoint.
\item At a terminal PE slot whose recorded support, Gamma color, and
defective color are \(v,\gamma,\delta\), the
\(\gamma\)-\(\delta\) component through \(v\), under every proper
exactly-one-uncolored coloring stable on the terminal ladder tree relative
to the terminal stored set and coloring, meets that ladder tree precisely
at \(v\).
\item Under every such stable terminal-ladder coloring, if \(\delta\) is
defective and \(v\) is a \(\delta\)-boundary vertex having an earlier
\(\delta\)-boundary vertex, then for every occurrence of \(v\) in the
terminal ETT and every index \(a<m\) whose slot \(a+1\) records \(v_a\) as
a PE support or an SE extension vertex, first appearance of \(v\) no later
than \(a+1\) implies \(v=v_a\) or \(v\) precedes \(v_a\) in the exact
terminal tree order.
\item Every proper exactly-one-uncolored coloring stable on the terminal
ladder tree relative to its terminal stored set and coloring is a
\noderef{ModColoring} of \(\mathcal L\).
\end{enumerate}

Assume also assertions (iii)--(vi), with exactly the four scopes in items
3--6, uniformly for every maximum ETT on the same \(G,k,e_0\), with
arbitrary certified input and generating colorings, whose rung number is
exactly \(n\).  In particular, the predecessor-ladder stability conclusion
in assertion (iii) remains conditional on the first-appearance index being
positive; none of the arbitrary-input quantifiers is restricted to the
input of \(\mathcal E\).

Finally assume the large branch
\[
                         11\leq \noderef{TashkinovOrder}(G),
\]
and assume the following candidate-uniform estimate.  For every same-rung
maximum ETT \(\mathcal L\) admitted above, every hierarchy base on it, and
every good hierarchy \(K\) over that base, if \(L_n\) is its terminal
ladder tree and \(\psi_n\) its generating coloring, then
\[
             2n+11\leq
             |\noderef{TreeMissingColors}(\psi_n,L_n)|.       \tag{1}
\]
Then
\[
       \noderef{ColoredVertexSetElementary}
          (G,\varphi_n,\noderef{TreeSequenceVertexSet}(T))
\]
holds.
\end{theorem}
\begin{proof}
We prove the strengthened assertion simultaneously for every choice of
certified input coloring, generating coloring, same-rung maximum ETT,
hierarchy base, and good hierarchy on the fixed \(G,k,e_0\), by strong
induction on the divider count \(q\).  Thus, at the step \(q\), the induction
hypothesis says that every such candidate with fewer than \(q\) dividers has
an elementary terminal vertex set.  This is exactly the rung-local induction
(6.1); neither the rung \(n\) nor the graph, palette, or distinguished exact
edge changes.

Fix a candidate at divider count \(q\), and suppose for a contradiction that
its terminal vertex set is not elementary.  Positivity is part of
\noderef{ExtendedTashkinovTree}, while the hierarchy, goodness, and maximum
certificates are among the fixed data.  Hence these data form a
\noderef{HierarchyCounterexample} at divider count \(q\).

Among all hierarchy counterexamples with this same rung number and divider
count \(q\), allowing the certified input coloring, generating coloring,
ETT, base, and hierarchy to vary, choose one \(C\) with least first
coordinate of \noderef{HierarchyCounterexampleMeasure}; among those choose
one with least second coordinate.  These choices exist by the well-ordering
of the natural numbers, because the counterexample just exhibited makes the
class nonempty.  By the definitions of
\noderef{HierarchyCounterexampleMeasure} and
\noderef{HierarchyCounterexampleLess}, the choice says precisely that no
same-rung, equal-divider counterexample \(C'\) satisfies
\[
 \noderef{HierarchyCounterexampleLess}
  (\noderef{HierarchyCounterexampleMeasure}(C'),
   \noderef{HierarchyCounterexampleMeasure}(C)).              \tag{2}
\]
The strong induction hypothesis rules out every counterexample with fewer
than \(q\) dividers, so \(q\) is also the least divider count relevant at
this induction step.

We next construct, field by field, the
\noderef{HierarchyMinimumCounterexampleContext} for \(C\).  First, its
generating coloring is proper and leaves exactly \(e_0\) uncolored.  Indeed,
let \(\mathcal S\) be its certified \noderef{TashkinovSeries}.  Since the
rung number is positive, its last certificate in
\noderef{TashkinovSeriesStep} is a positive SE, PE, or RE constructor.  In
the SE case, the current-coloring equation identifies the generating
coloring with the selected coloring in the recorded
\noderef{MaximumDefectiveChoice}, which is proper and
exactly-one-uncolored.  In the PE case, these two properties are explicit
current-coloring fields of the constructor.  In the RE case, the current
coloring equals the preceding coloring, for which the same two properties
are constructor fields.  In all three cases the generating-coloring field
of \noderef{ExtendedTashkinovTree} makes the resulting coloring exactly the
displayed generating coloring of \(C\).

The complete lower-rung six-pack is the first uniform hypothesis, applied
to arbitrary input and generating data.  The fewer-divider field is exactly
the strong induction hypothesis.  Equation (2) is exactly the
equal-divider lexicographic-minimality field.  Finally apply the uniform
estimate (1) to the chosen ETT, its maximum certificate, its chosen base,
its actual hierarchy, and its good-hierarchy certificate.  This gives
\[
             2n+11\leq|\overline{\varphi_n}(T_n)|             \tag{3}
\]
for the terminal ladder tree of \(C\), rather than reusing an estimate from
the initially displayed ETT.  These are all fields of the minimum context.
The assumed current-rung assertions (iii)--(vi) also apply to \(C\), since
their quantifiers allow precisely its arbitrary input and generating
colorings and its rung number is \(n\).

For later use we derive the palette-surplus assertion required in Situations
7.1 and 7.3.  Inspect the initial certificate of \(\mathcal S\).  At index
zero, \noderef{TashkinovSeriesStep} can only be its initial constructor, so
it supplies a genuine \noderef{KTriple} on \(G,k,e_0\).  Consequently
\[
                       \noderef{MaximumDegree}(G)+1\leq k.    \tag{4}
\]
Let \(\mu\) be any proper \(k\)-edge-coloring leaving exactly \(e_0\)
uncolored, and let \(v\) be any vertex.  By
\noderef{IsProperEdgeColoring}, distinct colored edges incident with \(v\)
have distinct colors.  Hence the number of \noderef{PresentColors} at \(v\)
is at most \noderef{Degree}\((G,v)\), which is at most
\(\operatorname{MaximumDegree}(G)\).  If
\(\noderef{MissingColors}(G,\mu,v)\) were empty, every one of the \(k\)
colors would be present at \(v\), giving
\[
 k\leq\deg_G(v)\leq\operatorname{MaximumDegree}(G),
\]
contrary to (4).  Therefore
\[
 \noderef{MissingColors}(G,\mu,v)\ne\varnothing              \tag{5}
\]
for every such \(\mu\) and every \(v\).  Notice that (5) is derived from
the certified \(k\)-triple and properness; it is not an extra field of the
minimum context.

Write \(H_C=C.\mathit{hierarchy}\), let \(A\) be its last starred stage, and
put
\[
 \ell=\noderef{HierarchyFullTailLength}(C)=|T_C|-|A|,
 \qquad r=\noderef{HierarchyPathNumber}(H_C).
\]
Apply \noderef{HierarchyCounterexamplePenultimateElementary} with the
properness, exact one-uncoloredness, lower six-pack, strong-induction
fewer-divider field, (2), and (3).  It gives the exact penultimate full-tail
prefix assertion
\[
 V(T_C(y_{\ell-1}))\text{ is elementary under }\varphi_n,     \tag{6}
\]
which is fact (6.6), derived after the minimum context rather than inserted
into it.

Applying \noderef{HierarchyCounterexampleTailAtLeastTwo} to those same six
context fields gives
\[
                              2\leq\ell.                      \tag{7}
\]
Applying \noderef{HierarchyCounterexamplePathNumberBounds} gives
\[
                              1\leq r\leq\ell.                \tag{8}
\]
Here \(\ell\) is the number of exact full-tail additions, whereas \(r\) is
the least path-start index; (7)--(8) do not identify them.

It remains only the exhaustive trichotomy checked by (7)--(8).  If \(r=1\),
apply \noderef{HierarchySituationOneImpossible} to \(C\), its complete
minimum context, this equality, and the universal missing-color assertion
(5).  This is a contradiction.  If \(r=\ell\), apply
\noderef{HierarchySituationTwoImpossible} to the same \(C\), context, and
equality, again obtaining a contradiction.  In the remaining case,
(8), \(r\ne1\), and \(r\ne\ell\) give
\[
                              2\leq r<\ell.
\]
Now \noderef{HierarchySituationThreeImpossible}, with the minimum context,
these two inequalities, and (5), gives the final contradiction.  Thus no
divider-\(q\) counterexample exists.  Strong induction proves elementarity
for the originally supplied good hierarchy, as required.
\end{proof}
